% ---
% title: C++17 if 与 switch 初始化语句
% date: 2026-07-12
% author: Crystal-Sky
% tags: C++, C++17, 基础语法, if, switch
% summary: 介绍 C++17 中 if 与 switch 初始化语句的语法、变量作用域、结构化绑定组合及常见应用。
% slug: cpp-if-switch-initializers
% course: C++
% course_slug: cpp
% chapter: 基础语法
% chapter_order: 1
% lesson_order: 2
% lesson_type: 基础语法
% ---

\documentclass[UTF8,12pt,a4paper]{ctexart}

\usepackage{geometry}
\usepackage{xcolor}
\usepackage{listings}
\usepackage{amsmath}
\usepackage{booktabs}
\usepackage{enumitem}
\usepackage{hyperref}
\usepackage{tcolorbox}
\usepackage{fancyhdr}

\geometry{left=2.4cm,right=2.4cm,top=2.5cm,bottom=2.5cm}

\hypersetup{
    colorlinks=true,
    linkcolor=blue,
    urlcolor=blue
}

\pagestyle{fancy}
\fancyhf{}
\fancyhead[L]{C++17 if 与 switch 初始化语句}
\fancyhead[R]{教学讲义}
\fancyfoot[C]{\thepage}

\definecolor{codebg}{RGB}{247,248,250}
\definecolor{keywordcolor}{RGB}{0,92,197}
\definecolor{commentcolor}{RGB}{0,128,0}
\definecolor{stringcolor}{RGB}{163,21,21}

\lstset{
    language=C++,
    basicstyle=\ttfamily\small,
    keywordstyle=\color{keywordcolor}\bfseries,
    commentstyle=\color{commentcolor},
    stringstyle=\color{stringcolor},
    backgroundcolor=\color{codebg},
    frame=single,
    numbers=left,
    numberstyle=\tiny\color{gray},
    numbersep=8pt,
    tabsize=4,
    showstringspaces=false,
    breaklines=true,
    columns=fullflexible,
    keepspaces=true
}

\newtcolorbox{conceptbox}[1]{
    colback=blue!3,
    colframe=blue!60!black,
    title=#1,
    fonttitle=\bfseries
}

\newtcolorbox{warningbox}[1]{
    colback=red!3,
    colframe=red!65!black,
    title=#1,
    fonttitle=\bfseries
}

\newtcolorbox{examplebox}[1]{
    colback=green!3,
    colframe=green!50!black,
    title=#1,
    fonttitle=\bfseries
}

\title{\textbf{C++17 \texttt{if} 与 \texttt{switch} 初始化语句}}
\author{}
\date{}

\begin{document}

\maketitle

\section{学习目标}

完成本节学习后，应能够：

\begin{enumerate}[label=\arabic*.]
    \item 理解 C++17 中选择语句初始化器的作用；
    \item 掌握 \texttt{if} 初始化语句的基本语法；
    \item 掌握 \texttt{switch} 初始化语句的基本语法；
    \item 理解初始化变量的作用域；
    \item 在查找、加锁、解析返回值等场景中正确使用该语法；
    \item 避免常见的作用域和语法错误。
\end{enumerate}

\section{什么是选择语句初始化器}

在 C++17 之前，程序员通常需要先定义变量，再在 \texttt{if} 或 \texttt{switch} 中使用它。

例如：

\begin{lstlisting}
auto it = mp.find("apple");

if (it != mp.end()) {
    std::cout << it->second << '\n';
}
\end{lstlisting}

这里的变量 \texttt{it} 在整个外层作用域中都存在，即使它只在判断语句中使用。

C++17 允许在 \texttt{if} 和 \texttt{switch} 的条件前直接写一条初始化语句：

\begin{lstlisting}
if (auto it = mp.find("apple"); it != mp.end()) {
    std::cout << it->second << '\n';
}
\end{lstlisting}

\begin{conceptbox}{核心思想}
选择语句初始化器允许程序员：

\begin{itemize}
    \item 在进入判断前创建一个临时变量；
    \item 在条件表达式中使用该变量；
    \item 将变量的作用域限制在当前 \texttt{if} 或 \texttt{switch} 语句中。
\end{itemize}
\end{conceptbox}

\section{\texttt{if} 初始化语句}

\subsection{基本语法}

C++17 中 \texttt{if} 初始化语句的基本形式如下：

\begin{lstlisting}
if (初始化语句; 条件表达式) {
    // 条件成立时执行
}
\end{lstlisting}

带 \texttt{else} 的形式：

\begin{lstlisting}
if (初始化语句; 条件表达式) {
    // 条件成立
} else {
    // 条件不成立
}
\end{lstlisting}

其中，初始化语句和条件表达式之间使用分号分隔。

\begin{warningbox}{注意}
初始化语句后必须写分号。

正确：

\begin{lstlisting}
if (int x = getValue(); x > 0)
\end{lstlisting}

错误：

\begin{lstlisting}
if (int x = getValue() x > 0)
\end{lstlisting}
\end{warningbox}

\subsection{最简单的示例}

\begin{lstlisting}
#include <iostream>
using namespace std;

int getValue() {
    return 10;
}

int main() {
    if (int x = getValue(); x > 0) {
        cout << "x 是正数，x = " << x << '\n';
    }

    return 0;
}
\end{lstlisting}

在这个例子中：

\begin{enumerate}[label=\arabic*.]
    \item 先执行 \texttt{int x = getValue()}；
    \item 再判断 \texttt{x > 0}；
    \item 如果条件成立，则进入 \texttt{if} 语句块。
\end{enumerate}


\section{结合结构化绑定使用}

C++17 中，初始化语句还可以和结构化绑定一起使用。

例如，\texttt{map::insert} 会返回一个 \texttt{pair}：

\begin{itemize}
    \item 第一个元素是迭代器；
    \item 第二个元素表示是否插入成功。
\end{itemize}

\begin{lstlisting}
#include <iostream>
#include <map>
#include <string>
using namespace std;

int main() {
    map<string, int> scores;

    if (auto [it, inserted] =
            scores.insert({"Alice", 95});
        inserted) {
        cout << "插入成功" << '\n';
    } else {
        cout << "键已经存在" << '\n';
    }

    return 0;
}
\end{lstlisting}

在条件表达式中直接判断 \texttt{inserted}。

如果插入失败，仍可在 \texttt{else} 分支通过 \texttt{it} 访问已经存在的元素：

\begin{lstlisting}
if (auto [it, inserted] =
        scores.insert({"Alice", 95});
    inserted) {
    cout << "插入成功" << '\n';
} else {
    cout << "已有分数：" << it->second << '\n';
}
\end{lstlisting}

\section{处理函数返回值}

如果一个函数返回状态码，可以在 \texttt{if} 中完成调用、保存结果和判断。

\begin{lstlisting}
#include <iostream>
using namespace std;

int connectServer() {
    return 0;
}

int main() {
    if (int code = connectServer(); code == 0) {
        cout << "连接成功" << '\n';
    } else {
        cout << "连接失败，错误码：" << code << '\n';
    }

    return 0;
}
\end{lstlisting}

这里的 \texttt{code} 只用于当前判断，因此不需要放到外层作用域。


\section{\texttt{switch} 初始化语句}

\subsection{基本语法}

C++17 允许在 \texttt{switch} 中加入初始化语句：

\begin{lstlisting}
switch (初始化语句; 条件表达式) {
    case 常量1:
        // ...
        break;

    case 常量2:
        // ...
        break;

    default:
        // ...
        break;
}
\end{lstlisting}

\subsection{简单示例}

\begin{lstlisting}
#include <iostream>
using namespace std;

int getCommand() {
    return 2;
}

int main() {
    switch (int command = getCommand(); command) {
        case 1:
            cout << "添加" << '\n';
            break;

        case 2:
            cout << "删除" << '\n';
            break;

        case 3:
            cout << "查询" << '\n';
            break;

        default:
            cout << "未知命令" << '\n';
            break;
    }

    return 0;
}
\end{lstlisting}

这里先定义变量 \texttt{command}，然后以 \texttt{command} 的值进行分支选择。

\section{\texttt{switch} 初始化变量的作用域}

初始化变量的作用域覆盖：

\begin{itemize}
    \item \texttt{switch} 条件表达式；
    \item 所有 \texttt{case} 分支；
    \item \texttt{default} 分支。
\end{itemize}

离开整个 \texttt{switch} 后，变量失效。

\begin{lstlisting}
switch (int code = getCommand(); code) {
    case 1:
        cout << code << '\n';
        break;

    default:
        cout << code << '\n';
        break;
}

// cout << code;  // 错误：code 已离开作用域
\end{lstlisting}

\section{\texttt{switch} 的实际应用}

假设一个函数返回命令类型：

\begin{lstlisting}
enum class Command {
    Add,
    Remove,
    Search,
    Unknown
};

Command readCommand() {
    return Command::Search;
}
\end{lstlisting}

可以使用：

\begin{lstlisting}
switch (auto cmd = readCommand(); cmd) {
    case Command::Add:
        cout << "执行添加操作" << '\n';
        break;

    case Command::Remove:
        cout << "执行删除操作" << '\n';
        break;

    case Command::Search:
        cout << "执行查询操作" << '\n';
        break;

    default:
        cout << "未知操作" << '\n';
        break;
}
\end{lstlisting}

这种写法把命令获取和命令分支处理放在同一个结构中。

\section{\texttt{if} 与 \texttt{switch} 初始化语句对比}

\begin{center}
\begin{tabular}{lll}
\toprule
项目 & \texttt{if} 初始化语句 & \texttt{switch} 初始化语句 \\
\midrule
引入版本 & C++17 & C++17 \\
分隔符 & 分号 & 分号 \\
适用判断 & 任意布尔条件 & 整数、枚举等离散值 \\
变量作用域 & 整个 \texttt{if--else} & 整个 \texttt{switch} \\
常见用途 & 查找、插入、指针判断 & 命令码、状态码、枚举 \\
\bottomrule
\end{tabular}
\end{center}


\section{常见错误}

\subsection{忘记写分号}

错误：

\begin{lstlisting}
if (int x = getValue() x > 0) {
}
\end{lstlisting}

正确：

\begin{lstlisting}
if (int x = getValue(); x > 0) {
}
\end{lstlisting}

\subsection{在语句外使用变量}

\begin{lstlisting}
if (int x = 10; x > 0) {
    cout << x << '\n';
}

cout << x << '\n';  // 错误
\end{lstlisting}

变量 \texttt{x} 的作用域只到 \texttt{if} 语句结束。

\subsection{误以为只能定义变量}

初始化语句不一定必须是变量声明，也可以是普通表达式语句。

例如：

\begin{lstlisting}
int x = 0;

if (x = getValue(); x > 0) {
    cout << x << '\n';
}
\end{lstlisting}

不过教学和实际开发中，更常见的是在初始化语句中声明局部变量。

\subsection{在多个变量声明中使用不同类型}

下面的写法错误：

\begin{lstlisting}
if (int x = 1, double y = 2.0; x < y) {
}
\end{lstlisting}

同一条声明语句中不能这样混合声明不同类型。

可以改为使用结构体、\texttt{pair} 或 \texttt{tuple}，或者提前声明其中一个变量。


\section{综合示例：用户查询系统}

\begin{lstlisting}
#include <iostream>
#include <map>
#include <string>
using namespace std;

int main() {
    map<string, int> users{
        {"Alice", 1001},
        {"Bob", 1002},
        {"Tom", 1003}
    };

    string name;
    cin >> name;

    if (auto it = users.find(name);
        it != users.end()) {
        cout << "用户编号：" << it->second << '\n';
    } else {
        cout << "用户不存在" << '\n';
    }

    return 0;
}
\end{lstlisting}

这个程序中，迭代器 \texttt{it} 只用于查询结果判断，因此放在 \texttt{if} 初始化语句中最合适。


\section{知识总结}

\begin{conceptbox}{本节核心内容}
\begin{enumerate}[label=\arabic*.]
    \item C++17 允许在 \texttt{if} 和 \texttt{switch} 中加入初始化语句；
    \item 初始化语句和条件表达式之间使用分号分隔；
    \item 初始化变量的作用域覆盖整个选择语句；
    \item 离开选择语句后，初始化变量自动销毁；
    \item 该语法特别适合容器查找、插入结果、函数状态码和命令处理；
    \item 与结构化绑定结合时，可以写出更紧凑、清晰的代码。
\end{enumerate}
\end{conceptbox}

\section{板书式速记}

\begin{lstlisting}
// if 初始化语句
if (auto value = getValue(); value > 0) {
    // ...
} else {
    // ...
}

// 容器查找
if (auto it = mp.find(key); it != mp.end()) {
    // ...
}

// 结构化绑定
if (auto [it, inserted] = s.insert(x); inserted) {
    // ...
}

// switch 初始化语句
switch (auto cmd = getCommand(); cmd) {
    case 1:
        break;
    default:
        break;
}
\end{lstlisting}

\end{document}
